\documentclass[a4paper,11pt]{ltjsarticle}
\usepackage[haranoaji]{luatexja-preset}
\usepackage{amsmath,amssymb}
\usepackage[top=12mm,bottom=10mm,left=14mm,right=14mm]{geometry}
\usepackage{array}
\usepackage{tikz}
\usetikzlibrary{calc}
\pagestyle{empty}
\setlength{\parindent}{0pt}
\newlength{\qcol}\setlength{\qcol}{0.47\textwidth}
\newlength{\qgap}
\newlength{\anscolw}\setlength{\anscolw}{0.30\textwidth}
\newlength{\ansh}
\newcommand{\Q}[2]{\parbox[t]{\qcol}{\raggedright (#1)\hspace{0.6em}$\displaystyle #2$}}
\newcommand{\QR}[4]{\Q{#1}{#2}\hfill\Q{#3}{#4}\par\vspace{\qgap}}
\newcommand{\anscell}[1]{\rule[-\ansdrop]{0pt}{\ansh}(#1)}
\newlength{\ansdrop}

\begin{document}
{\large\bfseries 計算問題　28}\hfill 名前(\underline{\hspace{32mm}})\par\vspace{3mm}
■（1）〜（16）の計算をしなさい。（17）〜（21）は方程式を解きなさい。\par\vspace{3mm}
\setlength{\qgap}{5.4mm}
\QR{1}{-4^2+(-5)^2-6}{2}{\dfrac{x-4}{5}\times30}
\QR{3}{(5x-4)-(2x+6)}{4}{12-(-5)\div\left(-\dfrac{5}{6}\right)}
\QR{5}{(-8)\times(-7)}{6}{(20a-30)\div5}
\QR{7}{4\div5\times(-30)}{8}{5\times\{5-(-6+1)\times2\}}
\QR{9}{(4a+5)+(6a-2)}{10}{\dfrac{1}{4}(20x-24)+\dfrac{1}{2}(12x+10)}
\QR{11}{(-11)+(-8)-(-11)}{12}{4(5x-2)-(6x-2)}
\QR{13}{(-20)\div(+5)}{14}{\dfrac{4}{5}(5x-10)-\dfrac{6}{2}(2x+8)}
\QR{15}{-0.4-(-6.5)-1.2}{16}{\dfrac{4}{5}\div6+\dfrac{2}{5}}
\QR{17}{4x-5=3}{18}{4(x+5)=5x+16}
\QR{19}{0.4x+0.5=1.7}{20}{\dfrac{2}{3}(x-1)=\dfrac{1}{2}x+1}
\vspace{1mm}
\Q{21}{\dfrac{4x-5}{6}-\dfrac{x+2}{6}=3}\par\vspace{3mm}
\setlength{\ansh}{9.5mm}\setlength{\ansdrop}{6mm}%
\begin{center}\begin{tabular}{|p{\anscolw}|p{\anscolw}|p{\anscolw}|}\hline
\anscell{1} & \anscell{2} & \anscell{3} \\\hline
\anscell{4} & \anscell{5} & \anscell{6} \\\hline
\anscell{7} & \anscell{8} & \anscell{9} \\\hline
\anscell{10} & \anscell{11} & \anscell{12} \\\hline
\anscell{13} & \anscell{14} & \anscell{15} \\\hline
\anscell{16} & \anscell{17} & \anscell{18} \\\hline
\anscell{19} & \anscell{20} & \anscell{21} \\\hline
\end{tabular}\end{center}

\end{document}
